\documentclass[11pt,a4paper,UTF8]{ctexart}
\usepackage{geometry}
\geometry{left=18mm,right=18mm,top=24mm,bottom=24mm,headheight=22pt,footskip=12mm}
\usepackage{amsmath,amssymb,mathtools,bm,esint,extarrows,centernot}
\usepackage{xcolor}
\usepackage{tcolorbox}
\tcbuselibrary{skins,breakable}
\usepackage{fancyhdr}
\usepackage{titlesec}
\usepackage{hyperref}
\usepackage{tikz}
\usepackage{booktabs}
\usepackage{tabularx}

% --- Color Palette (Purple & DeepPink Academic Theme) ---
\definecolor{DeepPurple}{HTML}{4C1D95}
\definecolor{TitlePurple}{HTML}{581C87}
\definecolor{SubPurple}{HTML}{6B21A8}
\definecolor{DeepPink}{HTML}{FF1493}
\definecolor{LightPinkBg}{HTML}{FFF5F8}
\definecolor{LightPurpleBg}{HTML}{F8F5FF}
\definecolor{GoldAccent}{HTML}{D97706}
\definecolor{DarkText}{HTML}{1F2937}

\hypersetup{
    colorlinks=true,
    linkcolor=TitlePurple,
    citecolor=DeepPink,
    urlcolor=SubPurple,
    pdfborder={0 0 0}
}

% --- Typography & Spacing ---
\linespread{1.3}
\setlength{\parskip}{0.35em plus 0.1em minus 0.1em}
\setlength{\parindent}{0pt}

% --- Section Titles Styling ---
\titleformat{\section}
  {\Large\bfseries\color{TitlePurple}}{\thesection}{1em}{}
  [\vspace{1mm}{\color{TitlePurple!30}\hrule height 1pt}]
\titleformat{\subsection}
  {\large\bfseries\color{SubPurple}}{\thesubsection}{0.8em}{}
\titleformat{\subsubsection}
  {\normalsize\bfseries\color{DeepPurple}}{\thesubsubsection}{0.6em}{}

% --- tcolorbox Environments ---
% 1. Theorem / Formula / Method / Analysis Box (Deep Purple)
\newtcolorbox{thmbox}[1][]{
    enhanced,
    breakable,
    colback=LightPurpleBg,
    colframe=TitlePurple,
    arc=3pt,
    boxrule=0pt,
    leftrule=3.5pt,
    left=10pt,
    right=10pt,
    top=8pt,
    bottom=8pt,
    title=#1,
    coltitle=white,
    colbacktitle=TitlePurple,
    fonttitle=\bfseries\small,
    attach boxed title to top left={yshift=-2mm, xshift=4mm},
    boxed title style={boxrule=0pt, arc=2pt}
}

% 2. Solution / Formula / Derivation Box (DeepPink)
\newtcolorbox{solbox}[1][]{
    enhanced,
    breakable,
    colback=LightPinkBg,
    colframe=DeepPink,
    arc=3pt,
    boxrule=0pt,
    leftrule=3.5pt,
    left=10pt,
    right=10pt,
    top=8pt,
    bottom=8pt,
    title=#1,
    coltitle=white,
    colbacktitle=DeepPink,
    fonttitle=\bfseries\small,
    attach boxed title to top left={yshift=-2mm, xshift=4mm},
    boxed title style={boxrule=0pt, arc=2pt}
}

% 3. Learning Goals / Summary Box
\newtcolorbox{targetbox}[1][]{
    enhanced,
    breakable,
    colback=LightPurpleBg,
    colframe=DeepPurple,
    arc=3pt,
    boxrule=1pt,
    left=10pt,
    right=10pt,
    top=8pt,
    bottom=8pt,
    title=#1,
    coltitle=white,
    colbacktitle=DeepPurple,
    fonttitle=\bfseries\small,
    attach boxed title to top left={yshift=-2mm, xshift=4mm},
    boxed title style={boxrule=0pt, arc=2pt}
}

\begin{document}

% ------------------- 讲义封面 (Cover Page) -------------------
\begin{titlepage}
\thispagestyle{empty}
\begin{tikzpicture}[remember picture,overlay]
    \fill[DeepPurple] (current page.north west) rectangle ([yshift=-7cm]current page.north east);
    \draw[DeepPink, line width=3.5pt] ([yshift=-7cm]current page.north west) -- ([yshift=-7cm]current page.north east);
    \draw[GoldAccent, line width=1pt] ([yshift=-7.12cm]current page.north west) -- ([yshift=-7.12cm]current page.north east);
    
    \node[anchor=north east, opacity=0.12, text=white] at ([xshift=-1.5cm, yshift=-0.5cm]current page.north east) {
        \fontsize{70}{70}\selectfont\bfseries 2027
    };
\end{tikzpicture}

\vspace*{0.8cm}
\begin{center}
    {\color{white}\fontsize{26}{32}\selectfont\bfseries 2027 考研数学(二) 核心讲义}\\[0.6cm]
    {\color{white}\fontsize{13}{18}\selectfont\bfseries 全国硕士研究生招生考试 · 统考数学(二) 核心精讲全书}\\[1.5cm]
\end{center}

\vspace{3.5cm}
\begin{center}
    {\color{GoldAccent}\large\bfseries $\blacktriangleright$ 基础强化 · 核心题解 · 考点深水区 $\blacktriangleleft$}\\[0.8cm]
    {\color{TitlePurple}\fontsize{24}{28}\selectfont\bfseries 第六章 一元函数微分学的应用(中值定理\&微分等式\&微分不等式)}\\[0.6cm]
    {\color{DeepPink}\large\bfseries —— 体系化理论精讲与公式全景 ——}\\[1.5cm]
\end{center}

\vfill

\begin{center}
\begin{tcolorbox}[
    colback=white,
    colframe=DeepPurple,
    width=0.88\textwidth,
    arc=4pt,
    boxrule=1.2pt,
    halign=center,
    drop shadow=gray!30
]
    \vspace{3mm}
    {\bfseries\large 编著团队：EscoffierZhou}\\[2.5mm]
    {\bfseries\color{TitlePurple} 目标院校：中国科学院大学 (UCAS) 计算机专硕 (22408)}\\[2.5mm]
    {\small\color{gray} 备考铁律：手算真题数字 · 错题白纸重做 · 408 客观题为王}\\[2.5mm]
    {\small 编制版本：2027 届首发版 · \today}
    \vspace{2mm}
\end{tcolorbox}
\end{center}
\vspace{0.8cm}
\end{titlepage}

% ------------------- 目录与页面主体配置 -------------------
\pagestyle{fancy}
\fancyhf{}
\fancyhead[L]{\small\color{gray}2027 考研数学(二) · 核心讲义系列}
\fancyhead[R]{\small\color{TitlePurple}\nouppercase{\leftmark}}
\fancyfoot[C]{\small\color{gray}— 第 \thepage\ 页 —}
\fancyfoot[R]{\small\color{gray}UCAS 22408}
\renewcommand{\headrulewidth}{0.5pt}

\pagenumbering{Roman}
\tableofcontents
\newpage
\pagenumbering{arabic}


\begin{targetbox}[本章复习目标与核心知识图谱]
\begin{itemize}
  \item 目的[1]:掌握闭区间连续函数的零阶中值定理(最值、介值、平均值、零点定理)及其几何意义
\end{itemize}
\end{targetbox}

\begin{targetbox}[本章复习目标与核心知识图谱]
\begin{itemize}
  \item 目的[2]:掌握费马引理、达布定理、罗尔、拉格朗日、柯西、泰勒中值定理的条件、结论与几何本质
\end{itemize}
\end{targetbox}

\begin{targetbox}[本章复习目标与核心知识图谱]
\begin{itemize}
  \item 目的[3]:熟练掌握辅助函数构造法则(积分因子法、微分方程还原法)及中值等式的证明思路
\end{itemize}
\end{targetbox}

\begin{targetbox}[本章复习目标与核心知识图谱]
\begin{itemize}
  \item 目的[4]:掌握方程根的存在性、唯一性与根的个数判定(零点定理、严格单调性、罗尔推论)
\end{itemize}
\end{targetbox}

\begin{targetbox}[本章复习目标与核心知识图谱]
\begin{itemize}
  \item 目的[5]:掌握利用单调性、凹凸性、最值及泰勒展开证明微分不等式的核心方法
\end{itemize}
\end{targetbox}


\section{1.涉及函数的中值定理(零阶中值定理)}

\textit{}<font color=red>设$f(x)$在闭区间$[a, b]$上连续:</font>\textit{}

\begin{thmbox}[定理[1]: 有界与最值定理]

定理:连续函数在闭区间$[a, b]$上必有界，且必能取得最大值$M$和最小值$m$:
\begin{equation*}
m \le f(x) \le M \quad(x \in [a, b])
\end{equation*}


  开区间或间断点不保证有界及最值(例如$y = \frac{1}{x}$在$(0, 1)$上无最值且无界) 

\end{thmbox}
\begin{thmbox}[定理[2]: 介值定理]

定理:设$f(x)$在$[a, b]$上连续，且$m, M$分别为其最小值与最大值 

对于任意介于$m$与$M$之间的实数$\mu$(即$m \le \mu \le M$)，至少存在一点$\xi \in [a, b]$，使得:
\begin{equation*}
f(\xi) = \mu
\end{equation*}


\end{thmbox}
\begin{thmbox}[定理[3]: 平均值定理(离散与连续)]

离散形式:当($a < x_1 < x_2 < \dots < x_n < b$)时，在$[x_1, x_n]$内至少存在一点$\xi$，使得:
\begin{equation*}
f(\xi) = \frac{f(x_1) + f(x_2) + \dots + f(x_n)}{n}
\end{equation*}


连续形式(积分第一中值定理):存在$\xi \in [a, b]$，使得:
\begin{equation*}
f(\xi) = \frac{1}{b - a} \int_a^b f(x)\,\mathrm{d}x
\end{equation*}


\end{thmbox}
\begin{thmbox}[定理[4]: 零点定理及其推广]

标准形式:若$f(x)$在$[a, b]$上连续，且$f(a) \cdot f(b) < 0$，则在开区间$(a, b)$内至少存在一点$\xi$，使得:
\begin{equation*}
f(\xi) = 0
\end{equation*}


推广形式(极限保号零点定理):

[前提]设$f(x)$在$(a, b)$内连续，若$\lim_{x \to a^+} f(x) = \alpha$，$\lim_{x \to b^-} f(x) = \beta$($\alpha, \beta$可为有限数或$\pm\infty$)，且$\alpha \cdot \beta < 0$ 

[结论]则由极限局部保号性，在开区间$(a, b)$内至少存在一点$\xi$，使得$f(\xi) = 0$ 

\end{thmbox}

\section{2.涉及导数(微分)的中值定理}

\begin{thmbox}[定理[1]: 费马引理(导数极值必要条件)]

定理:设$f(x)$在点$x_0$处满足:[可导]and[在$x_0$处取得极值],则必有:
\begin{equation*}
f'(x_0) = 0
\end{equation*}


  作用:内部极值点必定是驻点 证明导数等于零的中值定理(如达布定理)常以“最值在内部取得 + 费马引理”为破题主线 

\end{thmbox}
\begin{thmbox}[定理[2]: 达布定理(导数零点定理/导数介值定理)]

定理:设$f(x)$在$[a, b]$上可导(不要求导函数$f'(x)$连续):

  若$f'_+(a) \cdot f'_-(b) < 0$，则在开区间$(a, b)$内至少存在一点$\xi$，使得:
\begin{equation*}
f'(\xi) = 0
\end{equation*}

  
推论:导函数虽然不一定连续，但一定具备介值性(导函数不能有第一类间断点) 

\end{thmbox}
\begin{thmbox}[定理[3]: 罗尔中值定理]

定理:设$f(x)$满足以下所有条件:

  条件[1]:在$[a, b]$上连续

  条件[2]:在$(a, b)$内可导

  条件[3]:$f(a) = f(b)$ 

则在开区间$(a, b)$内至少存在一点$\xi$，使得:
\begin{equation*}
f'(\xi) = 0
\end{equation*}


推广形式:设$f(x)$在$(a, b)$内可导，若$\lim_{x \to a^+} f(x) = \lim_{x \to b^-} f(x) = A$($A$可为有限数或$\pm\infty$)，

则在$(a, b)$内至少存在一点$\xi$，使得$f'(\xi) = 0$ (区间可推广至$(-\infty, +\infty)$) 

\textcolor{red}{额外注意(130+):}

\end{thmbox}
\begin{thmbox}[定理[4]: 拉格朗日中值定理]

定理:设$f(x)$满足以下所有条件: 

  条件[1]:在$[a, b]$上连续

  条件[2]:在$(a, b)$内可导

则在开区间$(a, b)$内至少存在一点$\xi$，使得:
\begin{equation*}
f(b) - f(a) = f'(\xi)(b - a) \Leftrightarrow \left(f'(\xi) = \frac{f(b) - f(a)}{b - a}\right)
\end{equation*}


有限增量形式:设$x \in U(x_0)$，记$\Delta x = x - x_0$，则存在$\theta \in(0, 1)$，使得:
\begin{equation*}
\Delta y = f(x_0 + \Delta x) - f(x_0) = f'(x_0 + \theta \Delta x)\Delta x
\end{equation*}


  本质作用:建立函数增量$f(b) - f(a)$与导数值$f'(\xi)$之间的等价转换桥梁(以导数控制函数增量) 

\end{thmbox}
\begin{thmbox}[定理[5]: 柯西中值定理]

定理:设$f(x), g(x)$满足以下所有条件:

条件[1]:在$[a, b]$上连续

条件[2]:在$(a, b)$内可导

条件[3]:对任意$x \in(a, b)$均有$g'(x) \neq 0$ 

则在开区间$(a, b)$内至少存在一点$\xi$，使得:
\begin{equation*}
\frac{f(b) - f(a)}{g(b) - g(a)} = \frac{f'(\xi)}{g'(\xi)}
\end{equation*}


  几何本质:平面参数曲线$\begin{cases} X = g(t) \\ Y = f(t) \end{cases}$上的拉格朗日中值定理，连接端点的割线斜率等于某切线斜率 

\end{thmbox}
\begin{thmbox}[定理[6]: 泰勒中值定理]

定理:设$f(x)$在点$x_0$的某个邻域内具有直到$n+1$阶导数，则对该邻域内任意$x$，有:

\end{thmbox}

\begin{equation*}
f(x) = f(x_0) + f'(x_0)(x - x_0) + \frac{f''(x_0)}{2!}(x - x_0)^2 + \dots + \frac{f^{(n)}(x_0)}{n!}(x - x_0)^n + R_n(x)
\end{equation*}



拉格朗日余项(区间全局/定量证明专用):


\begin{equation*}
R_n(x) = \frac{f^{(n+1)}(\xi)}{(n+1)!}(x - x_0)^{n+1} \quad(\xi \text{ 介于 } x_0 \text{ 与 } x \text{ 之间})
\end{equation*}

佩亚诺余项(局部逼近/求极限专用):

\begin{equation*}
R_n(x) = o\left((x - x_0)^n\right) \quad(x \to x_0)
\end{equation*}


\section{3.中值定理证明的核心思维与辅助函数构造}

\begin{thmbox}[原理[1]: 中值定理题型识别路径]

\textbf{[1]} 不含导数:		   优先考虑[最值定理]or[介值定理]or[零点定理]

\textbf{[2]} 含一阶导数与一个中值$\xi$:	优先考虑[罗尔定理(需构造辅助函数)]or[拉格朗日中值定理]

\textbf{[3]} 含两不同函数同一中值$\xi$:

情况A:对于[纯分式比值]and[两函数自变量区间一致]:比如$\frac{f(b)-f(a)}{g(b)-g(a)} = \frac{f'(\xi)}{g'(\xi)}$

  写成分式比值形式，优先考虑[柯西中值定理] 

情况B:两个导数是[相加][相减][交叉相乘][带系数混合]:比如$f'(\xi) + g'(\xi) = 0$

利用微分方程还原法合并为一个一元辅助函数，只用一次罗尔定理

\textbf{[4]} 含两个不同中值$\xi, \eta$:	   考虑分段使用[拉格朗日中值定理]，或分别对两端点使用[拉格朗日中值定理]

\textbf{[5]} 含高阶导数$f^{(n)}(\xi)(n \ge 2)$:     条件含导数点值时优先考虑[泰勒公式展开] 条件含多个函数零点时优先考虑多次使用[罗尔定理]

\end{thmbox}
\begin{thmbox}[原理[2]: 罗尔辅助函数的标准构造表(微分方程还原法)]

| \textbf{待证结论形式} | \textbf{逆向积分来源(微分方程分离变量)} | \textit{}构造辅助函数 $F(x)$\textit{} |
| :--- | :--- | :--- |
| $f'(\xi) + \varphi'(\xi)f(\xi) = 0$ | 乘积分因子 $e^{\varphi(x)}$ | $F(x) = f(x)e^{\varphi(x)}$ |
| $f'(\xi) + kf(\xi) = 0$ | $\frac{f'}{f} + k = 0 \implies \ln\vert{}f\vert{} + kx = C$ | $F(x) = f(x)e^{kx}$ |
| $f'(\xi) - f(\xi) = 0$ | $\frac{f'}{f} - 1 = 0 \implies \ln\vert{}f\vert{} - x = C$ | $F(x) = f(x)e^{-x}$ |
| $xf'(\xi) + nf(\xi) = 0$ | 同乘 $x^{n-1} \implies(x^n f)' = 0$ | $F(x) = x^n f(x)$ |
| $xf'(\xi) - nf(\xi) = 0$ | 同除 $x^{n+1} \implies \left(\frac{f}{x^n}\right)' = 0$ | $F(x) = \frac{f(x)}{x^n}$ |
| $f'(\xi)g(\xi) + f(\xi)g'(\xi) = 0$ | 乘积求导法则 $[f(x)g(x)]'$ | $F(x) = f(x)g(x)$ |
| $f'(\xi)g(\xi) - f(\xi)g'(\xi) = 0$ | 商求导法则分子 $\left[\frac{f(x)}{g(x)}\right]'$ | $F(x) = \frac{f(x)}{g(x)}$ |
| $f''(\xi)f(\xi) - [f'(\xi)]^2 = 0$ | 商导数 $\left(\frac{f'}{f}\right)'$ 或 $(\ln f)''$ | $F(x) = \frac{f'(x)}{f(x)}$ 或 $F(x) = \ln\vert{}f(x)\vert{}$ |
| $f(\xi)f'(\xi) > 0$(不等式) | 复合平方导数 $\left[\frac{1}{2}f^2(x)\right]'$ | $F(x) = \frac{1}{2}f^2(x)$ |

\end{thmbox}

\section{4.微分等式(方程根的存在性与个数)}

\begin{thmbox}[原理[1]: 方程根与函数零点的对应法则]

方程$f(x) = 0$的根即为函数$y = f(x)$的零点 方程$f(x) = g(x)$的根即为两曲线$y = f(x)$与$y = g(x)$的交点横坐标 

\end{thmbox}
\begin{thmbox}[原理[2]: 实根存在性与唯一性判定工具]

存在性(至少有一个根):

零点定理:若$f(x)$在$[a, b]$上连续且$f(a) \cdot f(b) < 0$，则在$(a, b)$内至少有一个实根 

实系数奇次多项式定理:奇次多项式$x^{2n+1} + a_1 x^{2n} + \dots + a_{2n+1} = 0$在$(-\infty, +\infty)$内\textbf{至少有一个实根} 

唯一性(至多有一个根):

  严格单调性:若$f(x)$在区间$I$上严格单调增加或单调减少，则$f(x) = 0$在$I$内\textbf{至多有一个实根} 

  判定闭环:“区间连续端点异号 + 导数保号(严格单调)” $\implies$ 方程在该区间内\textbf{有且仅有一个实根} 

\end{thmbox}
\begin{thmbox}[原理[3]: 罗尔原话与根的上限判定(数根神器)]

罗尔定理推论:若在区间$I$上$f^{(n)}(x) \neq 0$(即方程$f^{(n)}(x) = 0$无实根)，则方程$f(x) = 0$在该区间内\textit{}至多有$n$个实根\textit{} 

罗尔原话(广义形式):若方程$f^{(n)}(x) = 0$在区间$I$上至多有$k$个实根，则方程$f(x) = 0$在该区间内\textit{}至多有$k + n$个实根\textit{} 

*(反证法逻辑:若$f(x) = 0$有$k+n+1$个根，由罗尔定理逐阶求导，$f^{(n)}(x) = 0$至少产生$k+1$个根，发生矛盾)* 

\end{thmbox}
\begin{thmbox}[原理[4]: 参变量分离与动静相交法]

对于含参数方程$f(x, k) = 0$，优先将参数$k$分离成$k = g(x)$形式:

方程实根的个数转化为水平直线$y = k$与定函数曲线$y = g(x)$的交点个数，通过研究$g(x)$的极值、渐近线与单调区间直接画图数根 

\end{thmbox}

\section{5.微分不等式证明方法与技巧}

\begin{thmbox}[原理[1]: 单调性法(移项构造一次求导)]

若证明$f(x) > g(x)(x > a)$，构造辅助函数$h(x) = f(x) - g(x)$:
若$h(a) \ge 0$且当$x > a$时$h'(x) > 0$，则$h(x)$在$[a, +\infty)$上严格单增，故$h(x) > h(a) \ge 0$ 

\end{thmbox}
\begin{thmbox}[原理[2]: 二阶导数保号与凹凸性几何割线法]

若在$(a, b)$内$f''(x) < 0$(严格凸弧)，且$f(a) = f(b) = 0$，则区间内部曲线恒位于连接两端点割线上方，恒有$f(x) > 0$ 

若在$(a, b)$内$f''(x) > 0$(严格凹弧)，且$f(a) = f(b) = 0$，则区间内部曲线恒位于割线下方，恒有$f(x) < 0$ 

\end{thmbox}
\begin{thmbox}[原理[3]: 拉格朗日中值定理放缩法]

对于形如$f(b) - f(a)$的差值结构，使用拉格朗日中值定理写成$f'(\xi)(b - a)$:

由$a < \xi < b$及导函数$f'(x)$的单调性，对$f'(\xi)$进行上下界放缩，直接得到双边不等式 

\end{thmbox}
\begin{thmbox}[原理[4]: 泰勒展开法(局部凸性切线放缩)]

由一阶泰勒展开式$f(x) = f(x_0) + f'(x_0)(x - x_0) + \frac{f''(\xi)}{2}(x - x_0)^2$:

若$f''(x) \ge 0$，则$f(x) \ge f(x_0) + f'(x_0)(x - x_0)$(凹函数曲线恒在切线上方) 

若$f''(x) \le 0$，则$f(x) \le f(x_0) + f'(x_0)(x - x_0)$(凸函数曲线恒在切线下方) 

\end{thmbox}
\begin{thmbox}[原理[5]: 考研核心经典微分不等式汇总]

\textbf{[1]} 三角函数基本链:当$x \in \left(0, \frac{\pi}{2}\right)$时，$\sin x < x < \tan x$

\textbf{[2]} 若尔当(Jordan)不等式:当$x \in \left(0, \frac{\pi}{2}\right)$时，$\frac{2}{\pi}x < \sin x < x$

\textbf{[3]} 对数不等式链:当$x > 0$时，$\frac{x}{1+x} < \ln(1+x) < x$

\textbf{[4]} 指数不等式:对任意实数$x$，恒有$e^x \ge x + 1$(等号仅在$x = 0$成立)
\end{thmbox}

\end{document}
